% C1: generated by c1_make_table.py from c1_prop5_certificate.json
\begin{table}[tb]
\centering\small
\caption{Rational interval certificate for Proposition~\ref{prop:TV} at $t_*=1047/100$ (uniformisation rate $2$, $101$ terms, $x=1047/50$, unnormalised tail bound $T<3.6\times10^{-27}$ against $B>10^{9}$). Shown are the digits common to both exact rational endpoints; the full endpoints are in \texttt{c1\_prop5\_certificate.json}. Signs refer to $p_i(t_*)-\pi_i$ and are identical for both preparations.}
\label{tab:prop5cert}
\begin{tabular}{cccccc}
\toprule
$i$ & $\pi_i$ & $p_{3,i}(t_*)$ & $p_{4,i}(t_*)$ & sign & width\\\midrule
$1$ & $0.7786380141$ & $0.779203311941127030\ldots$ & $0.779204375423765157\ldots$ & $+$ & $5.01432\times10^{-36}$\\
$2$ & $0.1816822033$ & $0.181790248374586092\ldots$ & $0.181789381981970039\ldots$ & $+$ & $3.33063\times10^{-36}$\\
$3$ & $0.0259546004$ & $0.025954537253576197\ldots$ & $0.025954598217632042\ldots$ & $-$ & $2.89144\times10^{-36}$\\
$4$ & $0.0129773002$ & $0.012977558197202530\ldots$ & $0.012977300930351277\ldots$ & $+$ & $2.85487\times10^{-36}$\\
$5$ & $0.0007478818$ & $0.000074344233508148\ldots$ & $0.000074343446281482\ldots$ & $-$ & $2.81850\times10^{-36}$\\
\midrule
\multicolumn{6}{p{.95\linewidth}}{$6.0176829178524738114767074365\times10^{-8}<\TV(p_3(t_*),\pi)-\TV(p_4(t_*),\pi)<6.0176829178524738114767074387\times10^{-8}$ (exact rational bounds)}\\
\multicolumn{6}{p{.95\linewidth}}{Certified signs of $\TV(p_3(t),\pi)-\TV(p_4(t),\pi)$: $t=10.46$: $-$;\ $t=10.47$: $+$;\ $t=10.48$: $-$.}\\
\multicolumn{6}{p{.95\linewidth}}{60-digit scan on $[0,60]$: crossings at $t_1=10.4608729171576\ldots$, $t_2=10.4752766252163\ldots$; no others.}\\
\bottomrule
\end{tabular}
\end{table}

% Suggested text for Prop. 5 (C1):
The reversal window is short and the effect is small. Rational certificates at $t=10.46$, $10.47$ and $10.48$ give the signs $-,+,-$ for $\TV(p_3(t),\pi)-\TV(p_4(t),\pi)$, so both crossings lie in explicit rational brackets, $t_1\in(10.46,10.47)$ and $t_2\in(10.47,10.48)$. A 60-digit scan locates them at $t_1\approx10.460872$ and $t_2\approx10.475276$; the window has length $0.01440$. The largest reversed difference is $6.02\times10^{-8}$, about $10^{-4}$ of either distance ($\TV\approx6.7\times10^{-4}$ at $t_*$). The initial conventional gap is $9.63\times10^{-3}$. Because $e^{tG}=\Pi+e^{-\eta t}(e^{tB}-\Pi)$, both distances carry the common factor $e^{-\eta t}$. The crossing times and window are therefore exactly independent of $\eta>0$, and the reset only rescales the effect. Under perturbations $w_5\mapsto w_5\pm10^{-6}$ and $w_5\pm10^{-5}$ (relative size up to $1.7\times10^{-4}$), combined with $\eta\mapsto\eta\pm10^{-5}$, every perturbed model keeps a reversal window inside $(10.4606,10.4755)$. Each such model also has a rational certificate of the reversed sign at $t_*$. Cone preservation is unaffected, since the quotient changes only by $-\eta I$. The preparations $\beta=3,4$ are colder than the bath ($\beta_b=1$) and undergo heating. The transformed quotient in Eq.~\eqref{eq:tvH} is Metzler only on the boundary: several off-diagonal entries are exactly zero.
